\section{Universos, categorías, equivalencia de categorías}
\label{VI.1}

Para evitar ciertas dificultades lógicas, admitiremos aquí la noción de
\emph{universo}, que es un conjunto «suficientemente grande» para no salir de
él mediante las operaciones habituales de la teoría de conjuntos; un
«\emph{axioma de los universos}» garantiza que todo objeto pertenece a algún
universo. Para más detalles, véase un libro en preparación de C\ptbl Chevalley
y el conferenciante\footnote{%
\ifthenelse{\boolean{orig}}
{Los autores definitivos son C\ptbl Chevalley y P\ptbl Gabriel. El libro debe
publicarse en el año 2000. \Cf también SGA~4 VI.7.1.}
{Los autores definitivos son C\ptbl Chevalley y P\ptbl Gabriel. El libro debe
publicarse en el año 3000. Mientras tanto, \cf también SGA~4 I.}}. Así, el
símbolo~$\Ens$
\label{indnot:fb}\oldindexnot{$\Ens$|hyperpage}%
designa, no la categoría de todos los conjuntos (noción que carece de sentido),
sino la categoría de los conjuntos que pertenecen a un universo dado (que aquí
no haremos constar en la notación). Del mismo modo, $\Cat$
\label{indnot:fc}\oldindexnot{$\Cat$|hyperpage}%
designará la categoría de las categorías pertenecientes al universo en
cuestión; los «morfismos» de un objeto~$X$ de~$\Cat$ en otro~$Y$ son, por
definición, los \emph{funtores} de~$X$ en~$Y$.

Si
\marginpar{147}
$\cal{C}$ es una categoría, denotamos por $\Ob(\cal{C})$
\label{indnot:fd}\oldindexnot{$\Ob(\cal{C})$|hyperpage}%
\emph{el conjunto de objetos} de~$\cal{C}$, y por~$\Fl(\cal{C})$
\label{indnot:fe}\oldindexnot{$\Fl(\cal{C})$|hyperpage}%
\emph{el conjunto de flechas} de~$\cal{C}$ (o morfismos de~$\cal{C}$).
Escribiremos, por tanto, $X\in\Ob(\cal{C})$, evitando el abuso de notación
habitual $X\in\cal{C}$. Si $\cal{C}$ y $\cal{C}'$ son dos categorías, un
\emph{funtor} de~$\cal{C}$ en~$\cal{C}'$ será siempre lo que comúnmente se
llama un funtor \emph{covariante} de~$\cal{C}$ en~$\cal{C}'$; su dato incluye
el de la
\ifthenelse{\boolean{orig}}{categoría de llegada y la categoría de partida}
{categoría de partida y la categoría de llegada}, $\cal{C}$ y $\cal{C}'$.
Los funtores de~$\cal{C}$ en~$\cal{C}'$ forman un conjunto, denotado
$\Hom(\cal{C},\cal{C}')$, que es el conjunto de objetos de una categoría
denotada $\SheafHom(\cal{C},\cal{C}')$.
\label{indnot:ff}\oldindexnot{$\SheafHom(\cal{C},\cal{C}')$|hyperpage}%
Por definición, un \emph{funtor contravariante} de~$\cal{C}$ en~$\cal{C}'$ es
un funtor de la \emph{categoría opuesta} $\cal{C}^\circ$
\label{indnot:fg}\oldindexnot{$\cal{C}^\circ$|hyperpage}%
de~$\cal{C}$ en~$\cal{C}'$.

Admitiremos las nociones de \emph{límite proyectivo} y de \emph{límite
inductivo} de un funtor $F\colon \cal{I}\to \cal{C}$ y, en particular, los
casos particulares más habituales de estas nociones: productos cartesianos y
productos fibrados, las nociones duales de sumas directas y sumas amalgamadas,
y las propiedades formales habituales de estas operaciones.

Por ejemplo, en la categoría $\Cat$ introducida más arriba, existen los límites
proyectivos (relativos a categorías~$\cal{I}$ pertenecientes al universo
elegido); el conjunto de objetos (\resp el conjunto de flechas) de la categoría
límite proyectivo~$\cal{C}$ de las~$\cal{C}_i$ se obtiene tomando el límite
proyectivo de los conjuntos de objetos (\resp de los conjuntos de flechas) de
las categorías~$\cal{C}_i$. El caso más conocido es el producto de una familia
de categorías. En lo sucesivo utilizaremos constantemente el producto fibrado
de dos categorías sobre una tercera.

Para todo lo relativo a categorías y funtores, mientras aparece el libro en
preparación ya mencionado, véase \cite{VI.1} (que es necesariamente muy
incompleto, incluso en lo que respecta a las generalidades esbozadas en el
presente número).

Aprovechemos esta ocasión para hacer explícita la noción de equivalencia de
categorías,
\index{equivalencia de categorías|hyperpage}%
que no se expone de manera satisfactoria en~\cite{VI.1}. Se dice que un funtor
$F\colon\cal{C}\to \cal{C}'$ es \emph{fiel}
\index{fiel (funtor)|hyperpage}%
(\resp \emph{plenamente fiel})
\index{plenamente fiel (funtor)|hyperpage}%
si, para todo par de objetos~$S$, $T$ de~$\cal{C}$, la aplicación
$u\mto F(u)$ de~$\Hom(S,T)$ en~$\Hom(F(S),F(T))$ es inyectiva
(\resp biyectiva). Se dice que $F$ es una \emph{equivalencia} de categorías si
\marginpar{148}
$F$ es plenamente fiel y, además, todo objeto~$S'$ de~$\cal{C}'$ es isomorfo a
un objeto de la forma~$F(S)$. Se demuestra que equivale a decir que existe un
funtor~$G$ de~$\cal{C}'$ en~$\cal{C}$ \emph{cuasi-inverso de}~$F$, es decir,
\ifthenelse{\boolean{orig}}
{tal que $GF$ sea isomorfo a $\id_{\cal{C}}$}
{tal que $GF$ sea isomorfo a $\id_{\cal{C}}$ y $FG$ a
$\id_{\cal{C}'}$}. Cuando así ocurre, el dato de un funtor
$G\colon\cal{C}'\to \cal{C}$ y de un isomorfismo
$\varphi\colon \ifthenelse{\boolean{orig}}{GF}{FG}\to \id_{\cal{C}'}$
equivale al dato, para todo $S'\in\Ob(\cal{C}')$, de un par~$(S,u)$ formado por
un objeto~$S$ de~$\cal{C}$ y un isomorfismo $u\colon F(S)\to S'$, a saber,
\ifthenelse{\boolean{orig}}{$(G(S),\varphi(S))$}
{$(G(S'),\varphi(S'))$}. (Con estas notaciones, existe un funtor único
$\cal{C}'\to\cal{C}$ cuya aplicación en los objetos es la aplicación dada
\ifthenelse{\boolean{orig}}{$S\mto G(S)$}{$S'\mto G(S')$}, y tal que la
aplicación \ifthenelse{\boolean{orig}}{$S\mto \varphi(S)$}
{$S'\mto \varphi(S')$} sea una transformación natural de funtores
$FG\to\id_{\cal{C}'}$). Por último, si $G$ es un funtor cuasi-inverso de~$F$
y se eligen isomorfismos $\varphi\colon FG\isomto\id_{\cal{C}'}$ y
$\psi\colon GF\isomto\id_{\cal{C}}$, entonces las dos condiciones de
\ifthenelse{\boolean{orig}}{compatibilidades}{compatibilidad} sobre~$\varphi$,
$\psi$ enunciadas en \cite[I.1.2]{VI.1} son, de hecho, equivalentes entre sí;
y, para todo isomorfismo~$\varphi$ elegido, existe un único
isomorfismo~$\psi$ que satisface dichas condiciones.
